\section{Backtesting Performance Evaluation}

The trained models are evaluated on an out-of-sample test period covering
\texttt{2026-01-01} to \texttt{2026-03-20} using the constituents of the Dow Jones 30
index. This phase evaluates the generalization capability of the learned policies under
previously unseen market conditions.

Based on the training results presented in the previous section, the best-performing
configurations are selected for backtesting: DDPG (MLP Policy), PPO (Transformer Policy),
and PPO (MLP--Transformer Policy). These models achieved the strongest overall training
performance and financial metrics among their respective policy families.

Evaluation is performed using the same financial metrics monitored during training:
\textit{total return}, \textit{portfolio value}, and \textit{Sharpe ratio}. Total return
measures cumulative portfolio growth, portfolio value tracks capital evolution over time,
while Sharpe ratio evaluates risk-adjusted performance. Table~\ref{tab:backtest_results}
summarizes the final total return values, while
Figures~\ref{fig:backtest_total_return},
\ref{fig:backtest_portfolio_value}, and
\ref{fig:backtest_sharpe_ratio} illustrate the evolution of the monitored metrics during
backtesting.

\begin{table}[h]
\centering
\begin{tabular}{lccc}
\hline
\textbf{Model} & \textbf{Smoothed Value} & \textbf{Final Total Return} & \textbf{Steps} \\
\hline
DDPG (MLP) & 0.0186 & 0.0112 & 51 \\
PPO (MLP--Transformer) & -0.0948 & -0.0992 & 51 \\
PPO (Transformer) & -0.1344 & -0.1481 & 51 \\
\hline
\end{tabular}
\caption{Backtesting performance measured using total return.}
\label{tab:backtest_results}
\end{table}

\begin{figure}
    \centering
    \includegraphics[width=1\linewidth]{figures/backtest_total_return.png}
    \caption{Backtesting total return evolution across evaluated models.}
    \label{fig:backtest_total_return}
\end{figure}

\begin{figure}
    \centering
    \includegraphics[width=1\linewidth]{figures/backtest_portfolio_value.png}
    \caption{Portfolio value evolution during backtesting.}
    \label{fig:backtest_portfolio_value}
\end{figure}

\begin{figure}
    \centering
    \includegraphics[width=1\linewidth]{figures/backtest_sharpe_ratio.png}
    \caption{Sharpe ratio evolution during backtesting.}
    \label{fig:backtest_sharpe_ratio}
\end{figure}

The two Transformer-based PPO models exhibit a gradual degradation in performance during
the evaluation period. As shown in
Figures~\ref{fig:backtest_total_return},
\ref{fig:backtest_portfolio_value}, and
\ref{fig:backtest_sharpe_ratio}, both policies show decreasing trends across all
financial metrics. The hybrid MLP--Transformer model exhibits a slower decline compared
to the pure Transformer architecture, indicating improved robustness through the
separation of market and portfolio representations, although both models ultimately
produce negative returns.

In contrast, DDPG (MLP) demonstrates the strongest out-of-sample performance, consistent
with the observations from training where it achieved the highest overall financial
metrics. As shown in Figure~\ref{fig:backtest_total_return}, total return increases
during the first part of evaluation, reaching a maximum around timestep 35, before
gradually declining toward the end of the testing period. Nevertheless, the model
maintains positive performance with a final total return of approximately $1.12\%$
(Table~\ref{tab:backtest_results}). The same behavior is reflected in portfolio value and
Sharpe ratio, indicating consistent profitability despite the late-stage decline.

Overall, the backtesting results remain aligned with the training observations. DDPG
(MLP), which achieved the strongest training performance, also demonstrates the best
generalization capability and out-of-sample robustness. PPO Transformer variants remain
competitive during training but exhibit reduced generalization performance under unseen
market conditions.